\subsection{Dyskryminacyjność cech i obsługa próbek trudnych w podejściu opartym na marginesach kątowych}

\paragraph{}Alternatywnym podejściem do problemu okluzji jest potraktowanie go nie jako błędu wymagającego naprawy (rekonstrukcji), lecz jako specyficznego przypadku "trudnej próbki" (ang. hard sample), która generuje dużą wariancję wewnątrzklasową. Wiodące badania w tym zakresie koncentrują się na modyfikacji funkcji strat w taki sposób, aby wymusić większą separowalność klas w przestrzeni cech, co pozwala modelowi zachować skuteczność nawet przy niekompletnych danych wejściowych.

Deng i in. [3] w swojej pracy nad metodą ArcFace (Additive Angular Margin Loss) wykazują, że standardowa funkcja straty Softmax nie jest wystarczająca do radzenia sobie z dużymi zmianami wyglądu twarzy (w tym okluzją), ponieważ nie optymalizuje bezpośrednio podobieństwa wewnątrzklasowego. Wprowadzenie addytywnego marginesu kątowego pozwala na rzutowanie cech na hipersferę, co znacząco zwiększa dyskryminacyjność modelu. Co jednak kluczowe dla problemu okluzji, autorzy zauważają, że zaszumione lub przysłonięte twarze mogą zaburzać proces uczenia, jeśli wymusza się ich dopasowanie do pojedynczego centrum klasy.

\paragraph{}Jako rozwiązanie tego problemu Deng i in. proponują mechanizm Sub-center ArcFace. W tym podejściu każda tożsamość może posiadać $K$ pod-centrów (ang. sub-centers), a próbka ucząca musi znajdować się blisko tylko jednego z nich. Badania wykazały, że mechanizm ten prowadzi do automatycznej separacji danych: "czyste" i wyraźne twarze grupują się wokół dominującego pod-centrum, podczas gdy twarze przysłonięte, ustawione pod dużym kątem lub zaszumione, trafiają do pod-centrów niedominujących [3].

Z perspektywy implementacji systemów odpornych na okluzję, wnioski płynące z analizy ArcFace sugerują, że:

1. Zwiększanie marginesu geometrycznego w funkcji straty jest skuteczną metodą poprawy odporności na brakujące cechy.

2. Okluzja może być modelowana jako osobna modalność w ramach tej samej tożsamości (poprzez sub-centers), co zapobiega degradacji głównego wzorca twarzy podczas treningu sieci.